\section{Algoritmos}
\begin{frame}[fragile,label=main]
	\begin{block}{}
		Sequência finita de instruções para executar uma tarefa.
		\begin{itemize}
			\item Bem definidas e não ambíguas.
			\item Executáveis com uma quantidade de esforço finita.
			\item Executáveis em um período de tempo finito.
		\end{itemize}
	\end{block}
	\vfill\pause
	\begin{columns}[T]\small
		\column{.3\textwidth}
			\begin{lstlisting}[title=Soma Sequencial]
				total = 0
				para i /in (0, 100]
				    total := total + i
			\end{lstlisting}
		\pause
		\column{.4\textwidth}
			\begin{lstlisting}[title=Bolo,morekeywords={tempo,batedeira,untar,esperar,conteudo}]
				tigela:= ingredientes
				enquanto(tempo() < 5min)
				    batedeira(tigela)
				untar(forma)
				forma:= conteudo(tigela)
				esperar(40min)
			\end{lstlisting}
	\end{columns}
\end{frame}

\begin{frame}[label=graduacao]
	Edsger Wybe \alert{Dijkstra}
	\begin{itemize}
		\item Holanda, $^{\star}$1930 $^{\dagger}$2002
		\item Prêmio Turing de 1972 (ACM)
			\begin{itemize}
				\item notação polonesa inversa  ($(a+b)/c \rightarrow a b + c /$)
				\item algoritmo de Dijkstra
				\item sistema operacional THE
				\item construção de semáforos para coordenar múltiplos processadores e programas
				\item ...
			\end{itemize}
	\end{itemize}
	\pause
\end{frame}

\newcommand{\function}[1]{\texttt{\textbf{#1}}}%
\begin{frame}[fragile,label=main]

	\begin{block}{Definição de Dijkstra}
		Um algoritmo corresponde a uma descrição de um padrão de comportamento, expresso em termos de um conjunto finito de ações.
		%\subitem[Ex:]{comportamento de $a + b$ para diferentes valores de $a$ e $b$.}
	\end{block}
	\vfill\pause	
	\begin{itemize}
		\item Um algoritmo não representa, necessariamente, um programa de computador.
			% \pause
			% \begin{itemize}
			%	\item Algoritmo: sequência de passos para resolver uma tarefa.
			%	\item Programa: sequência de instruções \alert<1>{em código} para executar uma operação em um computador.
			% \end{itemize}
			\pause
		\item Um algoritmo corretamente executado pode não resolver uma tarefa.% se:
			% \begin{itemize}
			%	\item foi incorretamente implementado;
			%	\item não for apropriado a tarefa.
			% \end{itemize}
	\end{itemize}
	\vfill
	Problema: aprovação em \insertshortsubtitle\pause
	\begin{columns}[T]\small
		\column{.35\textwidth}
			\begin{exampleblock}{Adequado}
				\function{assistir}(aula)\\
				\function{estudar}(algoritmos)\\
				\function{fazer}(provas $\land$ trabalhos)
			\end{exampleblock}
			\pause
		\column{0\textwidth}
		\column{.25\textwidth}
			\begin{block}{Indequado}
				\function{assistir}(filmes)\\
				\function{estudar}(química)\\
				\function{fazer}(faxina)
			\end{block}
			\pause
		\column{0\textwidth}
		\column{.25\textwidth}
			\begin{alertblock}{Não apropriado}
				\function{escovar}(dentes)\\
				\function{vestir}(pijamas)\\
				\function{deitar}
			\end{alertblock}
	\end{columns}
\end{frame}